\documentclass[11pt]{article}

\usepackage{amsmath,amssymb,graphicx}
\usepackage[margin=1in]{geometry}
\usepackage{booktabs}

\title{Emergent Relational Geometry Probe on a Frozen Kernel-Induced Cochain Operator Architecture}
\author{Tomislav Rupi\'c \\ Independent Researcher}
\date{March 2026}

\begin{document}

\maketitle

\begin{abstract}
Stages 10--14 established a frozen pre-geometric diagnostic program on a kernel-induced cochain operator architecture, covering single-packet morphology, collective interaction structure, coarse-grained envelope organization, scale-coupled fragility, and falsified weak binding extensions. Stage 15 changes viewpoint rather than dynamics: instead of adding new mechanisms, it tests whether frozen collective data already support a stable relational ordering that behaves metric-like under refinement. A minimal pilot is run on three representative regimes at a \((12 \rightarrow 24)\) resolution ladder: a compact composite cluster, a counter-propagating corridor, and a diffuse symmetric triad. Three induced relational distances are tested using only previously frozen observables: persistence distance, phase-correlation distance, and transport-delay distance. The resulting split is structured rather than uniform. Localized metric-like ordering appears in the counter-propagating corridor and in the phase-correlation channel of the diffuse triad, while the composite anchor remains weak or metric-violating and transport-delay ordering fails most consistently. The main conclusion is bounded: some proto-geometric compressibility is present, but it is regime-dependent and not anchored to the strongest composite envelopes. Stage 15 therefore does not support a robust emergent geometric layer on the frozen architecture.
\end{abstract}

\section{Introduction}

The program through Stage 14 establishes a cumulative diagnostic arc on a frozen kernel-induced cochain operator architecture. Stage 10 introduced a morphology instrument for packet transport. Stage 11 extended the same frozen packet family to collective interaction morphology. Stage 12 coarse-grained collective configurations into basin and envelope summaries. Stage 13 coupled refinement directly to interaction and coarse scales and showed that coarse organization sharpens under refinement while composite interaction labels remain fragile. Stage 14 then tested three weak structural extensions and found no effective binding channel at the tested strengths.

These results motivate a change of question. If weak structural extensions do not stabilize composite persistence, it becomes natural to stop asking what additional dynamics should be inserted and instead ask what ordering structure can already be read out from the frozen data. Stage 15 is the first probe of that kind.

The aim is narrow. Stage 15 does not construct spacetime. It does not assume a metric or a continuum limit. It asks only whether frozen collective interaction data induce a scale-coherent relational ordering that is stable enough to behave like a metric bookkeeping layer. Geometry, in this stage, is treated only as a possible compression of interaction persistence, phase ordering, and propagation order.

\section{Frozen Context and Pilot Cases}

All operators, projection rules, kernels, and packet families remain frozen. No new interaction term is introduced. The projected transverse sector and periodic boundary condition are retained, together with the packet family inherited from the strongest coherent transverse seed: bandwidth \(w=0.1\), cosine carrier, central wave number \(k=1.0\), and amplitude scale \(0.5\).

The first executable pilot uses three cases:
\begin{itemize}
\item asymmetric density cluster,
\item counter-propagating corridor,
\item symmetric wide triad.
\end{itemize}

The asymmetric density cluster is used as the clustered composite anchor because it is a genuine triadic configuration and therefore supports non-degenerate triangle and loop diagnostics. The counter-propagating corridor serves as a transport-dominated control. The symmetric wide triad serves as a weakly structured diffuse control.

Each case is evaluated at resolutions \(12\) and \(24\). This \((12 \rightarrow 24)\) ladder is chosen deliberately rather than reusing the earlier \((8 \rightarrow 16)\) ladder because Stage 13 showed that the lower-resolution side of the earlier ladder was still dominated by coarse-readability effects.

\section{Relational Distance Candidates}

Three induced relational distances are tested. All are constructed only from previously frozen time-series observables.

\subsection{Persistence distance}

The persistence distance is defined from the time-averaged pairwise overlap,
\[
 d_p(i,j) = \frac{1}{\epsilon + \langle \mathrm{overlap}_{ij} \rangle_T},
\]
with a small regularization \(\epsilon\). This distance is symmetric by construction and asks whether long-lived overlap persistence can induce a stable closeness ordering.

\subsection{Phase-correlation distance}

The phase-correlation distance is defined from the time-averaged absolute phase-lock indicator,
\[
 d_{\phi}(i,j) = 1 - \bigl|\langle \mathrm{phase\_lock}_{ij} \rangle_T\bigr|.
\]
This distance asks whether synchronization structure provides a more coherent relational ordering than overlap persistence.

\subsection{Transport-delay distance}

The transport-delay distance is defined from the first arrival time of a measurable envelope response at packet \(j\) after excitation at packet \(i\),
\[
 d_{\tau}(i \rightarrow j) = t_{\mathrm{arrival}}(i \rightarrow j).
\]
A symmetrized form is used for metric-like tests, but the raw directional version is also recorded. This channel asks whether propagation ordering can serve as a cone-like relational structure.

\section{Metric-Like Tests}

Each candidate distance is evaluated using three diagnostic tests.
\begin{itemize}
\item Symmetry deviation: mean and maximum asymmetry between \(d(i,j)\) and \(d(j,i)\).
\item Triangle-inequality violation rate: computed on all non-degenerate triplets.
\item Ordering persistence: rank-order stability of pairwise distances across the refinement ladder.
\end{itemize}

When applicable, two higher-level probes are added.
\begin{itemize}
\item Effective dimensionality from relational ball-growth, using persistence distance.
\item Loop-defect and geodesic-deviation analogs for triadic cases.
\end{itemize}

The pilot assigns only four descriptive labels:
\begin{itemize}
\item no coherent relational ordering,
\item weak ordering with high metric violation,
\item scale-fragile metric-like ordering,
\item scale-stable metric-like ordering.
\end{itemize}

No stronger label is allowed in this stage.

\section{Pilot Outcome}

The pilot produces the following aggregate counts:
\begin{itemize}
\item scale-stable metric-like ordering: 4,
\item weak ordering with high metric violation: 3,
\item no coherent relational ordering: 2.
\end{itemize}

This split is already informative. Stage 15 is not a uniform null, but neither is it a broad positive emergence signal. Instead, the induced ordering quality is strongly regime- and observable-dependent.

\subsection{Asymmetric density cluster}

The composite anchor produces only weak or failed ordering. Persistence distance and phase-correlation distance both enter the ``weak ordering with high metric violation'' class. Their ordering-flip fractions are \(0.3333\), their pair-rank correlation is \(0.5\), and the phase-correlation channel also shows a triangle-violation rate of \(0.3333\). The transport-delay channel fails most strongly here, entering ``no coherent relational ordering'' with an ordering-flip fraction of \(0.6667\) and symmetry deviation above \(0.32\).

The clustered composite case is therefore the key negative result of the pilot. The strongest composite anchor does not convert cleanly into a stable relational manifold under the tested induced distances.

\subsection{Counter-propagating corridor}

The counter-propagating corridor produces the cleanest local ordering. All three candidate distances are classified as scale-stable metric-like ordering. In each case the ordering-flip fraction is \(0\), the pair-rank correlation is \(1\), and symmetry deviation is numerically zero at the level recorded here.

This result must be interpreted carefully. The corridor case is only a two-packet system, so some metric tests are degenerate. Still, the diagnostic signal is real: stable ordering appears most clearly in the transport-like corridor, not in the strongest composite anchor.

\subsection{Symmetric wide triad}

The diffuse triad splits by channel. Persistence distance remains weak and refinement-sensitive, with ordering-flip fraction \(0.3333\), pair-rank correlation \(0.5\), and a large effective-dimension drift of about \(0.737\). By contrast, phase-correlation distance is classified as scale-stable metric-like ordering. Transport-delay distance fails again, entering ``no coherent relational ordering'' with symmetry deviation around \(0.328\).

This is the strongest triadic signal in the pilot: the phase-correlation channel provides the clearest non-degenerate large-scale ordering tendency.

\section{Interpretation of the Split}

The Stage 15 split is structurally more informative than a simple positive or negative outcome.

First, metric-like ordering can appear. The corridor case and the phase-ordered diffuse triad both support stable ordering under at least one candidate distance.

Second, the stable ordering is not anchored to the strongest composite regime. The asymmetric density cluster remains weak or metric-violating under all three distance candidates.

Third, transport-delay is the least successful distance in the present pilot. It fails on the composite anchor and on the diffuse triad, and only appears stable in the two-packet corridor control.

Taken together, these points imply a nontrivial separation:
\begin{itemize}
\item ordering can appear without strong composite persistence,
\item composite persistence can appear without stable metric-like ordering.
\end{itemize}

That separation is precisely what Stage 15 was designed to test. Persistent interaction structure is not automatically compressible into a geometric ordering layer.

\section{Interpretation Boundary}

Stage 15 remains diagnostic only. It does not support:
\begin{itemize}
\item spacetime claims,
\item physical metric claims,
\item causal-structure claims,
\item gravitational analogies.
\end{itemize}

The bounded claims supported by the pilot are:
\begin{itemize}
\item some localized metric-like ordering is present in the frozen data,
\item that ordering is strongly regime-dependent,
\item phase-correlation is the strongest large-scale ordering carrier in the present pilot,
\item persistence distance acts as a weaker structural carrier,
\item transport-delay is presently too fragile or discontinuous for geometry inference,
\item no robust scale-stable geometric layer appears on the composite anchor.
\end{itemize}

\section{Conclusion}

Stage 15 closes the first pre-geometric emergence boundary of the program. Instead of inserting a new interaction mechanism, it asks whether frozen collective dynamics already induce a stable relational ordering that behaves metric-like under refinement. The answer is mixed in a useful way.

Localized metric-like ordering does appear, but it is neither universal nor anchored to the strongest composite envelopes. The clearest ordering appears in a simple propagation corridor and in the phase-correlation channel of a symmetric triad. The compact composite anchor remains weak or metric-violating, and transport-delay remains the least reliable distance candidate.

This is therefore not the moment to claim emergent geometry. It is, however, a meaningful diagnostic boundary. The frozen architecture now shows proto-geometric compressibility in specific regimes without yet supporting a robust relational manifold. The next sensible step is not stronger geometric interpretation but refinement of relational transport and ordering diagnostics under perturbation and scale stress.

\end{document}
